% !TeX root = ../samplepaper.tex
\section{Proposed Method}
\subsection{Overview}
方法包含离线进化和在线控制两个过程，如图\ref{fig:overview}所示。离线阶段，GP在多个训练实例上评价控制器：每个individual参与若干次完整FACO求解，所得平均路径质量作为fitness。Selection与genetic operators据此产生下一代。训练结束后，使用validation选择一个individual，作为冻结控制器。

在线阶段，控制器在每次ACO迭代开始时读取当前搜索状态，并选择一个动作。该动作确定本次迭代的参考路径、起点区域和MNE，由整个蚁群共享；各蚂蚁仍独立采样起点和后续节点。构造与local search完成后，求解器更新最优路径、信息素和反馈统计，再进入下一次决策。因此，一棵GP树的输出首先是候选优先级，经过解码得到控制动作，完整求解结束后才能得到fitness。

\begin{figure}[tb]
\centering
\begin{tikzpicture}[x=1mm,y=1mm,>=Stealth,
 every node/.style={font=\small},
 box/.style={draw,align=center,text width=29mm,minimum height=11mm}]
\node[box] (pop) at (0,0) {GP population\\控制规则};
\node[box] (eval) at (40,0) {FACO evaluation\\平均终点fitness};
\node[box] (breed) at (80,0) {Selection\\Genetic operators};
\draw[->] (pop)--(eval); \draw[->] (eval)--(breed);
\draw[->] (breed.north)--++(0,5)-|node[pos=.25,above]{下一代}(pop.north);
\node[box] (state) at (0,-29) {搜索状态\\候选动作特征};
\node[box] (score) at (40,-29) {GP priority scores\\选择$(r,j,k)$};
\node[box] (aco) at (80,-29) {Construction\\Local search};
\draw[->] (state)--(score); \draw[->] (score)--(aco);
\draw[->] (aco.south)--++(0,-7)-|node[pos=.25,below]{更新路径、信息素与反馈}(state.south);
\draw[dashed,->] (eval.south)--node[right,font=\scriptsize]{一次评价中的迭代循环}(score.north);
\end{tikzpicture}
\caption{GP学习FACO控制策略。每个individual在训练实例上执行下方循环，终点质量形成fitness。冻结控制器使用相同的在线循环；参考标签仅由外部评价器读取。}
\label{fig:overview}
\end{figure}

\subsection{Search Decisions and Actions}
求解器维护活动参考路径、全局最优路径、阶段最优路径及容量为4的优质路径archive。Archive保留相对全局最优差距不超过2\%的不同路径。每次决策前，从中选择与活动路径采样边差异最大的合格路径作为备选，并为候选参考路径生成四个起点区域。

动作记为$a=(r,j,k)$。$r\in\{0,1\}$决定保留活动参考路径或切换至备选路径；$j\in\{0,1,2,3\}$决定起点区域；$k\in\{0,1,2,3\}$对应MNE为$2,4,8,16$。区域0和1为参考路径上独立随机抽取的连续片段，区域2沿近邻图广度优先扩展，区域3为不放回随机节点集；每个区域最多16个节点。区域只限制构造的起点，后续修改可以涉及其他节点。

每次迭代最多有32个完整动作。没有合格备选路径时，屏蔽$r=1$的动作。执行参考路径切换时，阶段最优改为备选路径，信息素重新初始化，阶段计数及回归率、local-search工作量的平滑统计清零；全局最优和全局停滞计数保留。当前控制接口没有重启冷却期。

对选定动作，每只蚂蚁复制其参考路径，从所选区域采样起点，完成聚焦构造后执行检查表驱动的2-opt。每只蚂蚁的local search最多评价$10^5$个移动。随后更新archive和反馈，以0.01概率采用阶段最优、否则采用本次迭代最优，作为信息素沉积路径和下一次活动参考路径。完整在线过程见算法\ref{alg:solve}。

\begin{algorithm}[tb]
\caption{FACO求解与GP控制}
\label{alg:solve}
\begin{algorithmic}[1]
\Require 实例$I$，控制器$\pi$，ACO seed $\xi$，蚂蚁数$m$，迭代数$H$
\State 初始化参考路径、信息素、最优路径与archive
\For{$t=1,\ldots,H$}
 \State 读取状态$S_t$，生成合法动作$\mathcal A(S_t)$及特征$\phi(S_t,a)$
 \State $a_t\gets\pi(S_t)$ \Comment{由GP分数解码得到$(r,j,k)$}
 \State 执行参考路径切换，确定起点区域和MNE
 \State 并行构造$m$条候选路径，分别执行local search
 \State 更新最优路径、archive、信息素及搜索反馈
\EndFor
\State \Return 全局最优路径
\end{algorithmic}
\end{algorithm}

\subsection{GP Representation and Decoding}
\paragraph{Joint single-tree（\JS{}）.}
一个individual包含一棵标量树$f$。对每个合法完整动作计算priority score，并选择最高分动作：
\begin{equation}
 \pi_J(S)=\arg\max_{a\in\mathcal A(S)}f\bigl(\phi(S,a)\bigr).
\end{equation}
同一棵树在不同候选特征上重复求值，因而能够比较三项决策的组合。每次决策最多需要32次树求值。

\paragraph{Conditional multi-tree（\CT{}）.}
一个individual包含角色固定的$f_r,f_j,f_k$三棵树，依次完成参考路径、区域和强度选择：
\begin{align}
 r^*&=\arg\max_{r\in\mathcal R(S)} f_r\bigl(\phi_r(S,r)\bigr),\\
 j^*&=\arg\max_{j\in\mathcal J(S,r^*)} f_j\bigl(\phi_j(S,r^*,j)\bigr),\\
 k^*&=\arg\max_{k\in\mathcal K(S,r^*,j^*)} f_k\bigl(\phi(S,r^*,j^*,k)\bigr).
\end{align}
各阶段的候选集合由合法完整动作投影得到；只要一个前缀存在合法补全，便可参与该阶段选择。三次选择读取同一个执行前状态，全部完成后才执行动作。$f_r$可见搜索状态及参考路径特征，$f_j$进一步读取区域特征，$f_k$读取全部输入。最多$2+4+4=10$次树求值即可得到动作。三棵树由同一个fitness共同评价，这与用multi-tree协同学习多个决策规则的GP框架一致\cite{zhang2018multitree,sun2024multicloud}。

\subsection{Terminal Set and Function Set}
表\ref{tab:features}列出12个归一化terminal及其可见性。搜索状态描述当前阶段的进展、停滞和计算反馈；参考路径及区域特征用于区分候选动作。回归率衡量local search后重新得到参考路径的蚂蚁比例，工作量反映搜索所消耗的局部移动评价数。这些反馈为GP提供选择扰动方式的依据。

\begin{table}[tb]
\caption{Terminal set。$\checkmark$表示该角色可以读取；\JS{}读取全部terminal。代码字段及精确归一化公式见补充材料。}
\label{tab:features}
\small
\begin{tabularx}{\textwidth}{@{}lYccc@{}}
\toprule
符号 & 含义 & $f_r$ & $f_j$ & $f_k$\\ \midrule
\multicolumn{5}{l}{搜索状态（search state）}\\
$p$ & Progress：已完成迭代比例$(t-1)/H$ & $\checkmark$ & $\checkmark$ & $\checkmark$\\
$s$ & Stagnation：未改善全局最优的迭代数，按32缩放 & $\checkmark$ & $\checkmark$ & $\checkmark$\\
$\rho$ & Return rate：返回参考路径的蚂蚁比例，平滑后使用 & $\checkmark$ & $\checkmark$ & $\checkmark$\\
$w$ & LS work：每蚂蚁平均局部移动评价数，按$10^5$缩放并平滑 & $\checkmark$ & $\checkmark$ & $\checkmark$\\
\midrule
\multicolumn{5}{l}{参考路径与动作（reference and action）}\\
$r$ & 候选动作的参考路径切换标记 & $\checkmark$ & $\checkmark$ & $\checkmark$\\
$\delta_q$ & 候选参考路径相对全局最优的差距，按0.02缩放 & $\checkmark$ & $\checkmark$ & $\checkmark$\\
$d_q$ & 候选与活动参考路径的采样边差异 & $\checkmark$ & $\checkmark$ & $\checkmark$\\
$u$ & MNE level：$(k+1)/4$，即0.25、0.5、0.75、1 & & & $\checkmark$\\
\midrule
\multicolumn{5}{l}{候选区域（candidate region）}\\
$e_R$ & Region excess：相邻边相对近邻距离的长度超额 & & $\checkmark$ & $\checkmark$\\
$a_R$ & Archive disagreement：区域相邻边在archive中的缺失比例 & & $\checkmark$ & $\checkmark$\\
$\tau_R$ & Pheromone strength：区域相邻边的归一化信息素 & & $\checkmark$ & $\checkmark$\\
$d_R$ & Region dispersion：相邻端点落在区域外的比例 & & $\checkmark$ & $\checkmark$\\
\bottomrule
\end{tabularx}
\end{table}

两种表示使用相同的function set：$+$、$-$、$\times$、$\min$、$\max$、绝对值及analytic quotient，后者定义为$\AQ(a,b)=a/\sqrt{1+b^2}$。常量terminal包括$-1,0,1$及从$[-2,2]$均匀采样的ephemeral random constants（ERC）。比例特征裁剪至$[0,1]$，两个平滑反馈采用权重$1/16$的指数移动平均。

实际树求值使用FP32，每个运算结果裁剪至$[-8,8]$。分数差不超过$10^{-6}$时，选择编号最小的合法候选；完整动作按$16r+4j+k$编号。这些规则确定近并列情况下的行为，解释表达式时同时保留其数学形式与实际解码结果。

\subsection{Fitness Evaluation and Evolution}
GP优化控制器在完整求解结束时的路径质量。设$\mathcal B_g$为第$g$代共享的实例--ACO seed面板，$T(I,\xi;\pi)$为算法\ref{alg:solve}返回的路径，$C_{\rm ref}(I)$为外部评价器读取的参考路径长度，则最小化
\begin{equation}
 F_g(\pi)=\frac{1}{|\mathcal B_g|}\sum_{(I,\xi)\in\mathcal B_g}
 100\left(\frac{C(T(I,\xi;\pi))}{C_{\rm ref}(I)}-1\right).
 \label{eq:fitness}
\end{equation}
该值以百分数表示，称为reference gap。参考路径未逐例独立认证最优性。参考标签不进入控制器输入或FACO求解。每代更换训练面板，同代所有individual及两种表示的运行共享面板，elites也重新评价。训练实例轮换在GP规则学习中已有应用\cite{mei2018orienteering,maclachlan2020}；本实验通过固定validation面板观察其学习表现。

初始化按目标高度2、3、4及full/grow方式分层采样；grow保留一条达到目标高度的分支。每棵树高度至多5，每个individual总节点数至多63，三树共享总节点限制。父本通过tournament selection选取。Crossover以0.9概率选择非terminal位置，并允许选择根；single-tree将第二父本的子树接入第一父本。Multi-tree先随机选择一个角色进行同角色subtree crossover，其余角色整树取自第二父本。该操作同时重组角色内部结构和角色间组合，与tree-swapping crossover属于相关设计\cite{zhang2018multitree}。

Mutation替换一个随机角色中的子树；对单叶树采用生长变异。每代保留elites，并以crossover、mutation或reproduction产生其余individual。为维持搜索多样性，接受的程序须互异，在64个固定训练情境上产生相同动作向量的individual最多两个。超出结构限制或重复的候选重新生成，必要时采用随机移民。具体概率、重试范围与精确语义列于实验设置及补充材料。

\begin{algorithm}[tb]
\caption{GP控制器的离线学习}
\label{alg:gp}
\begin{algorithmic}[1]
\Require 表示类型，GP seed，训练面板$\{\mathcal B_g\}_{g=1}^G$，validation面板
\State 初始化$P$个满足结构和行为约束的individual
\For{$g=1,\ldots,G$}
 \State 在$\mathcal B_g$上执行算法\ref{alg:solve}，按式\eqref{eq:fitness}评价整个种群
 \State 保存本代冠军，并按固定间隔在开发面板上监控
 \If{$g<G$}
  \State 保留elites，通过selection与genetic operators补齐下一代
 \EndIf
\EndFor
\State 对历代冠军与末代种群执行两级validation，选择并冻结一个控制器
\State \Return 冻结控制器
\end{algorithmic}
\end{algorithm}

\subsection{Parallel Fitness Evaluation}
Fitness evaluation将individual、实例／ACO seed和ant三个维度展开至GPU。不同蚁群分别维护路径、信息素、archive和随机流；实例几何量复用，GP树编译为紧凑指令执行。不同GP seed和独立评价任务分配至空闲RTX A5000。路径几何、代价和local-search增益使用FP64。

一次蚂蚁路径构造及其local search记为一次evaluation（FE），一次求解的预算为$mH$。此预算限制候选路径数；不同控制策略仍可能消耗不同数量的构造步和局部移动评价。因此，实验分别报告FE、搜索工作量及GPU时间。
